\chapter{Requirements and use cases}
\label{cha:req}

This chapter specifies what the laboratory must do. It identifies the actors, lists the functional and non-functional requirements with their origin and priority, and details the main use cases. The requirements were refined iteratively during development; the \emph{Source} column records why each one exists, so that it can be traced to the problem statement, to the threat analysis or to the evaluation design.

\section{Actors}
\label{sec:actors}

Three human actors and two external systems interact with the laboratory (Figure~\ref{fig:usecases}).

\begin{description}
\item[Learner] A student who uses the laboratory modules to generate and inspect phishing emails, examine deepfake scenarios and run adversarial attacks.
\item[Instructor] A teacher or examiner who demonstrates the modules, inspects the provenance of results and consults the exported evaluation files.
\item[Operator] The person who deploys the laboratory, sets the access credentials and API keys, and starts the attack service.
\item[Gemini API] The external LLM service used for phishing generation, email analysis and deepfake narration.
\item[Attack service] The local Python service that runs FGSM and PGD on the pretrained classifier.
\end{description}

\begin{figure}[htbp]
\centering
\begin{tikzpicture}[
  actor/.style={draw, rounded corners, font=\scriptsize\bfseries, minimum width=2.0cm, minimum height=0.7cm, align=center},
  uc/.style={draw, ellipse, font=\scriptsize, minimum width=3.6cm, minimum height=0.75cm, align=center},
  line/.style={thin}
]
\draw[rounded corners] (2.3,1.0) rectangle (8.9,-7.6);
\node[font=\scriptsize\bfseries] at (5.6,0.7) {InsightFlow AI};
\node[uc] (u1) at (5.6,0.0) {UC1 Sign in};
\node[uc] (u2) at (5.6,-1.1) {UC2 Load teaching set};
\node[uc] (u3) at (5.6,-2.2) {UC3 Generate phishing email};
\node[uc] (u4) at (5.6,-3.3) {UC4 Analyse email};
\node[uc] (u5) at (5.6,-4.4) {UC5 Inspect deepfake scenario};
\node[uc] (u6) at (5.6,-5.5) {UC6 Run adversarial attack};
\node[uc] (u7) at (5.6,-6.6) {UC7 Reproduce evaluation};
\node[actor] (l) at (0,-1.6) {Learner};
\node[actor] (i) at (0,-4.6) {Instructor};
\node[actor] (o) at (0,-6.6) {Operator};
\node[actor] (g) at (11.2,-2.7) {Gemini API};
\node[actor] (a) at (11.2,-5.5) {Attack service};
\foreach \u in {u1,u2,u3,u4,u5,u6} \draw[line] (l) -- (\u.west);
\foreach \u in {u1,u5,u6,u7} \draw[line] (i) -- (\u.west);
\draw[line] (o) -- (u7.west);
\draw[line] (u3.east) -- (g); \draw[line] (u4.east) -- (g); \draw[line] (u5.east) -- (g);
\draw[line] (u6.east) -- (a);
\end{tikzpicture}
\caption{Use-case diagram of the laboratory.}
\label{fig:usecases}
\end{figure}

\section{Functional requirements}
\label{sec:fr}

Table~\ref{tab:fr} lists the functional requirements, prioritised with the MoSCoW scheme (\emph{Must}, \emph{Should}, \emph{Could}). The source column refers to the problem statement (P), the threat analysis in Section~\ref{sec:stride} (T) or the evaluation design (E).

\begin{table}[htbp]
\centering
\caption{Functional requirements.}
\label{tab:fr}
\scriptsize
\begin{tabular}{P{0.7cm}P{6.4cm}P{1.0cm}P{0.7cm}P{2.7cm}}
\toprule
ID & Requirement & Priority & Source & Realised by \\
\midrule
FR1 & Require authentication for every API call when credentials are configured, and refuse to serve the laboratory in production when they are not. & Must & T & \texttt{labSecurity.ts} \\
FR2 & Generate a teaching phishing email from a scenario, a topic, optional instructions and a complexity level. & Must & P & \path{/api/generate-phishing} \\
FR3 & Analyse a subject, a body and an optional sender, and return scores, indicators, mitigations and a predicted class with a rationale. & Must & P & \path{/api/analyze-email} \\
FR4 & Accept a local \texttt{.txt} or \texttt{.eml} file and extract sender, recipient, subject and body in the browser. & Should & P & \texttt{parseEmailFile} \\
FR5 & Load a seeded, reproducible set of 90 teaching emails (30 per class) into the dashboard. & Must & P & \texttt{generate90Samples} \\
FR6 & Analyse a deepfake \emph{scenario} from a filename, a media type and a description, without reading media content. & Must & P & \path{/api/analyze-media-deepfake} \\
FR7 & Return fixed teaching content when the filename matches one of three presets. & Should & P & \texttt{PRE\_BAKED\_DEEPFAKES} \\
FR8 & Run FGSM or PGD on a pretrained CIFAR-10 classifier for a chosen class and budget, untargeted from the user interface and optionally targeted through the API. & Must & P, E & \path{/api/adversarial-perturb}, \texttt{service.py} \\
FR9 & Label every successful analytical response with \texttt{source} $\in$ \{\texttt{gemini}, \texttt{fallback}, \texttt{preset}, \texttt{computed}\}. & Must & P & \texttt{server.ts}, \texttt{service.py} \\
FR10 & Return a deterministic fallback when Gemini is unavailable on the email and deepfake routes; return HTTP~503 on the adversarial route when the attack service is unavailable. & Must & P & Route handlers \\
FR11 & Expose the health of the attack service and the exported metric files. & Should & E & \path{/api/eval/*} \\
FR12 & Present dashboard aggregates as configured teaching values, not as measurements. & Must & P & \path{UnifiedAnalysisDashboard.tsx} \\
FR13 & Provide scripts that regenerate the phishing benchmark and the adversarial batch results from code. & Must & E & \path{eval/} \\
\bottomrule
\end{tabular}
\end{table}

\section{Non-functional requirements}
\label{sec:nfr}

Table~\ref{tab:nfr} lists the non-functional requirements. Where possible, each has a measurable target that is verified in Chapter~\ref{cha:eval}.

\begin{table}[htbp]
\centering
\caption{Non-functional requirements with measurable targets.}
\label{tab:nfr}
\scriptsize
\begin{tabular}{P{0.9cm}P{2.1cm}P{6.6cm}P{0.8cm}}
\toprule
ID & Quality & Requirement and target & Source \\
\midrule
NFR1 & Security & JSON request bodies larger than 32\,kB are rejected. & T \\
NFR2 & Security & At most 10 \texttt{POST} requests per hour per client address to the generation route, and 40 to every other \texttt{POST /api/} route. & T \\
NFR3 & Security & Credentials are compared in constant time; empty or length-mismatched values never authenticate. & T \\
NFR4 & Security & Sessions expire after 12 hours; session tokens carry 192 bits of randomness and are sent in \texttt{HttpOnly}, \texttt{SameSite=Lax} cookies. & T \\
NFR5 & Security & Secrets are supplied through environment variables and are never committed to the repository. & T \\
NFR6 & Integrity of results & Every adversarial example satisfies $\|x_{\mathrm{adv}} - x\|_\infty \le \varepsilon$; a violation is reported as an error, never returned as a result. & E \\
NFR7 & Reproducibility & Running each evaluation script twice on the same commit yields identical metric files (fixed seed 20260927). & E \\
NFR8 & Availability & The website remains usable when Gemini or the attack service is down; affected routes degrade as specified in FR10. & P \\
NFR9 & Testability & The automated test suite runs offline, without calling Gemini, in under 10 seconds. & E \\
NFR10 & Portability & The application can be built and run as a container image. & P \\
NFR11 & Security & At most 10 failed sign-in attempts per client address within 15 minutes; further attempts receive HTTP~429. & T \\
\bottomrule
\end{tabular}
\end{table}

\section{Main use cases}
\label{sec:uc}

Tables~\ref{tab:uc1}--\ref{tab:uc6} detail the three use cases that carry the main security and evaluation logic. The remaining use cases follow the same pattern and are covered by the sequence diagrams in Chapter~\ref{cha:arch}.

\begin{table}[htbp]
\centering
\caption{UC1 --- Sign in.}
\label{tab:uc1}
\scriptsize
\begin{tabular}{P{2.4cm}P{10.0cm}}
\toprule
Actor & Learner or Instructor \\
Precondition & The Operator has set \texttt{LAB\_BASIC\_USER} and \texttt{LAB\_BASIC\_PASS}. \\
Main flow & (1) The browser loads the login page, which is served without a session. (2) The actor submits a username and password. (3) The server compares both values in constant time. (4) The server creates a session and returns an \texttt{HttpOnly} cookie. (5) The application shows the dashboard. \\
Alternative flows & (3a) Wrong credentials: HTTP~401, the form shows an error; after 10 failures from the same address within 15 minutes, HTTP~429. (1a) Credentials not configured in production: every request receives HTTP~503 and the laboratory stays closed. (1b) Non-interactive clients may use HTTP Basic authentication instead of the cookie. \\
Postcondition & Subsequent \path{/api/*} calls are accepted until the session expires or the actor signs out. \\
\bottomrule
\end{tabular}
\end{table}

\begin{table}[htbp]
\centering
\caption{UC4 --- Analyse email.}
\label{tab:uc4}
\scriptsize
\begin{tabular}{P{2.4cm}P{10.0cm}}
\toprule
Actor & Learner; Gemini API \\
Precondition & The actor is signed in. \\
Main flow & (1) The actor pastes a subject, body and optional sender, or loads a \texttt{.txt}/\texttt{.eml} file that is parsed in the browser. (2) The browser posts the fields to \path{/api/analyze-email}. (3) The server places the email between explicit untrusted-data markers in the prompt and calls Gemini with a JSON response schema. (4) The server normalises the predicted class and returns the analysis with \texttt{source=gemini}. (5) The dashboard adds the result to the session history. \\
Alternative flows & (3a) Gemini fails: the server returns a keyword-heuristic analysis with \texttt{source=fallback}. (2a) More than 40 requests in the last hour: HTTP~429. (2b) Body larger than 32\,kB: HTTP~413. \\
Postcondition & The result is visible with its provenance; nothing is stored on the server. \\
\bottomrule
\end{tabular}
\end{table}

\begin{table}[htbp]
\centering
\caption{UC6 --- Run adversarial attack.}
\label{tab:uc6}
\scriptsize
\begin{tabular}{P{2.4cm}P{10.0cm}}
\toprule
Actor & Learner or Instructor; Attack service \\
Precondition & The actor is signed in; the Operator has started the attack service. \\
Main flow & (1) The actor selects a CIFAR-10 class, FGSM or PGD, and a budget $\varepsilon \in (0, 0.25]$. (2) The server forwards the request to the attack service. (3) The service selects the first test image of that class that the model classifies correctly, runs an untargeted attack, and measures $\|x_{\mathrm{adv}} - x\|_\infty$. (4) The service returns the clean and adversarial predictions with their softmax confidences, the budget and the measured perturbation with \texttt{source=computed}. \\
Alternative flows & (2a) Attack service unreachable: HTTP~503 with \texttt{source=unavailable}; no values are shown. (3a) Unknown class or budget outside the allowed range: the service answers HTTP~400 and the server reports HTTP~502. (3b) Measured perturbation exceeds $\varepsilon$: the service answers HTTP~500 and the server reports HTTP~502; no result is returned. (1a) Through the API only, a request with \texttt{targeted: true} runs a targeted attack towards a second class and reports whether the target was reached. \\
Postcondition & The displayed numbers are model outputs for one image; they are distinct from the batch evaluation in Chapter~\ref{cha:eval}. \\
\bottomrule
\end{tabular}
\end{table}

\section{Out of scope}
\label{sec:oos}

The following are deliberately excluded: sending email, harvesting or storing real mailboxes, detecting manipulation in uploaded media, training or fine-tuning models, per-user accounts and roles, and server-side storage of generated or analysed content. Chapter~\ref{cha:disc} discusses the consequences of these exclusions.

\section{Summary}

The requirements combine teaching functions (FR2--FR8), a provenance contract (FR9, FR10, FR12), security controls derived from the threat analysis (FR1, NFR1--NFR5, NFR11) and evaluation properties (FR13, NFR6, NFR7). Chapter~\ref{cha:arch} shows how the architecture realises them.
